\section{Coordinate Systems}
\label{sec:coordinate-systems}

\subsection{Space}

Physics takes place in space.

To describe position we choose coordinates.

Nature itself has no coordinate axes.

Coordinates are introduced only to help us describe observations.

The most commonly used coordinate system is Cartesian coordinates.

A point is represented by

\[
(x,y,z).
\]

These numbers describe

\begin{itemize}
	\item distance along the x-axis,
	\item distance along the y-axis,
	\item distance along the z-axis.
\end{itemize}

%%%%%%%%%%%%%%%%%%%%%%%%%%%%%%%%%%%%%%%%%%%%%%%%%%%%%%

\subsection{Cartesian Coordinate System}

Define three mutually perpendicular axes.

\[
x,
\qquad
y,
\qquad
z.
\]

Their associated unit vectors are

\[
\hat{\mathbf x},
\qquad
\hat{\mathbf y},
\qquad
\hat{\mathbf z}.
\]

Each has unit length

\[
|\hat{\mathbf x}|=
|\hat{\mathbf y}|=
|\hat{\mathbf z}|=1.
\]

They satisfy

\[
\hat{\mathbf x}\cdot\hat{\mathbf y}=0,
\]

etc.

Therefore they are orthonormal.

\begin{bookfigure}{Rotating the coordinate axes changes the numerical components but not the geometric vector.}
\begin{tikzpicture}[x=1cm,y=1cm]
\draw[book axis] (-.3,0)--(5.4,0) node[right] {$x$};
\draw[book axis] (0,-.3)--(0,4.1) node[above] {$y$};
\begin{scope}[rotate=28]
\draw[book axis,draw=visualcyan] (-.2,0)--(5,0) node[right] {$x'$};
\draw[book axis,draw=visualcyan] (0,-.2)--(0,3.5) node[above] {$y'$};
\end{scope}
\draw[book vector] (0,0)--(4.05,3.15) node[above right] {$\mathbf v$};
\draw[book angle] (1,0) arc[start angle=0,end angle=28,radius=1];
\node[text=visualgold] at (1.23,.27) {$\phi$};
\node at (3.9,.55) {same vector, new components};
\end{tikzpicture}

\label{fig:coordinate-rotation}
\end{bookfigure}

%%%%%%%%%%%%%%%%%%%%%%%%%%%%%%%%%%%%%%%%%%%%%%%%%%%%%%
